% Teil II: durchgerechneter Weg, Teile a--d.

\subsection*{Gegeben und gesucht}
$\vec a=(1,3,-4)$, $\vec b=(3,5,-1)$. Gesucht: (a) Beträge und $\vec a+\vec b$,
(b) Skalarprodukt und Winkel, (c) Kreuzprodukt $\vec c$, (d) Nachweis
$\vec c\perp\vec a,\vec c\perp\vec b$ und geometrische Bedeutung von $|\vec c|$.

\subsection*{Ansatz}
Beträge und Summe komponentenweise; der Winkel folgt aus dem
Skalarprodukt, das Kreuzprodukt aus der Determinantenformel;
Senkrechtheit prüft man mit zwei Zeilen-Skalarprodukten.

\subsection*{Rechnung}
\textbf{a) Beträge und Summe.}
\begin{equation}
  |\vec a|=\sqrt{1+9+16}=\sqrt{26}\approx5{,}10,\qquad
  |\vec b|=\sqrt{9+25+1}=\sqrt{35}\approx5{,}92,
\end{equation}
\begin{equation}
  \vec a+\vec b=(4,8,-5)\;.
\end{equation}

\textbf{b) Skalarprodukt und Winkel.}
\begin{equation}
  \vec a\cdot\vec b=1\cdot3+3\cdot5+(-4)(-1)=3+15+4=22\;,
\end{equation}
\begin{equation}
  \cos\theta=\frac{22}{\sqrt{26}\,\sqrt{35}}
  \approx0{,}7293
  \;\Rightarrow\;
  \boxed{\;\theta\approx\qty{43{,}2}{\degree}\;}.
\end{equation}

\textbf{c) Kreuzprodukt.}
\begin{equation}
  \vec c=\vec a\times\vec b
  =\begin{pmatrix}
      3\cdot(-1)-(-4)\cdot5\\ (-4)\cdot3-1\cdot(-1)\\ 1\cdot5-3\cdot3
   \end{pmatrix}
  =\begin{pmatrix}17\\ -11\\ -4\end{pmatrix}\;.
\end{equation}

\textbf{d) Senkrechtheit und Bedeutung.}
\begin{align}
  \vec c\cdot\vec a &= 17\cdot1+(-11)\cdot3+(-4)(-4)=17-33+16=0\;,\\
  \vec c\cdot\vec b &= 17\cdot3+(-11)\cdot5+(-4)(-1)=51-55+4=0\;.
\end{align}
Da beide Skalarprodukte verschwinden, steht $\vec c$ senkrecht auf beiden
Vektoren.
\begin{equation}
  |\vec c|=\sqrt{17^2+(-11)^2+(-4)^2}=\sqrt{426}\approx\qty{20{,}6}{}\;.
\end{equation}

\begin{ergebnis}
$|\vec a|=\sqrt{26}\approx5{,}10$, $|\vec b|=\sqrt{35}\approx5{,}92$,
$\vec a+\vec b=(4,8,-5)$; Winkel $\theta\approx\qty{43{,}2}{\degree}$;
$\vec c=(17,-11,-4)$, $|\vec c|\approx20{,}6$.
\end{ergebnis}

\begin{kontrolle}
\textbf{Einheiten:} geometrische Grössen; Vektorkomponenten und Beträge
haben dieselbe Längeneinheit wie eingegeben. \checkmark

\textbf{Querprobe:} $|\vec c|=|\vec a|\,|\vec b|\,\sin\theta$
$=\sqrt{26}\sqrt{35}\sin(43{,}2^\circ)\approx\sqrt{910}\cdot0{,}684\approx20{,}6$. \checkmark

\textbf{Geometrie:} $|\vec c|$ ist der Flächeninhalt des von $\vec a$ und
$\vec b$ aufgespannten Parallelogramms.
\end{kontrolle}

\begin{merksatz}
Das Kreuzprodukt liefert einen Vektor senkrecht auf beiden Faktoren und
seinen Betrag die Parallelogrammfläche; das Skalarprodukt vermittelt den
Winkel zwischen zwei Vektoren.
\end{merksatz}
